\section{Language 作为统一认知接口}

\subsection{为什么 Language 仍是中心媒介}
课程指出，语言具备三重角色：状态编码介质、策略表达介质、协作协议介质。这使得语言不仅能承载自然语言解释，还能承载符号结构、程序片段与中间决策记录。

\begin{figure}[H]
\centering
\includegraphics[width=0.84\textwidth]{slide-03-reasoning-taxonomy.jpg}
\caption{视频约 00:37:20 讲义页：课程将推理类型按作用位置与目标进行分层}
\end{figure}

\begin{importantbox}{语言接口的系统优势}
\begin{itemize}
    \item 跨模态统一：文本可作为视觉、代码、结构化数据的交换层；
    \item 跨模块统一：planner、memory manager、tool router 可共享同一种协议；
    \item 跨代理统一：multi-agent 协作可以先通过语言协议快速原型化。
\end{itemize}
\end{importantbox}

\subsection{符号推理与语言推理的融合}
Yu Su 并未将 symbolic reasoning 与 neural reasoning 对立，而是强调 ``分工融合''：语言推理负责高维语义压缩与策略探索，符号推理负责精确约束与可验证计算。二者结合是长链任务稳定性的关键。

\begin{knowledgebox}{融合策略示例}
\begin{itemize}
    \item 语言模型生成候选计划，符号模块做约束检查；
    \item 语言模型做启发式搜索，程序执行器做终局验证；
    \item 语言模型做错误解释，规则系统做安全拦截。
\end{itemize}
\end{knowledgebox}

\begin{warningbox}{仅靠语言推理的边界}
当任务需要高精度数值、强约束组合优化或形式化证明时，单纯自然语言链条容易失真，需要引入程序化中间表示。
\end{warningbox}

\subsection{本章小结}
Language 不是 ``提示词输入格式''，而是 Agent 系统的统一通信层。其价值在于连接不同能力模块，而非替代所有模块。

